% Appendix, from exam/paper_b.md: Paper B, the questions and the register reference the candidate
% may consult. The model answers are a separate appendix, as they are a separate file, so the paper
% can be sat without them in view.

\chapter{Self-Assessment Paper B}
\label{exb:ch:paper}

\noindent\textbf{Time:} 3 hours. \textbf{Closed book.} No compiler and no datasheet. The register
reference at the back of this paper is the only material you may use. \textbf{Total: 100 marks.}

This paper exists to test your own skills and knowledge. It is not a qualification and it gates
nothing. It draws on all ten chapters, so \textbf{it is meant to be taken once you have finished
the book}. The model answers are in Appendix~\ref{ansb:ch:answers}.

\section*{Rubric}
\textbf{C11 on an ATmega328P (Arduino Uno/Nano) at 16\,MHz} is assumed throughout, built with
\code{avr-gcc} and \code{-std=c11}. Standard and AVR headers (\code{<stdint.h>},
\code{<stdbool.h>}, \code{<avr/io.h>}, \code{<avr/interrupt.h>}, \code{<util/delay.h>}) may be
assumed included unless a question asks you to name them.

\textbf{Register and bit names are not memory-tested.} Every register and bit this paper refers to
is listed in the register reference at the back (page~\pageref{exb:sec:registers}), with what it
is; consult it freely. What is marked is knowing which register to reach for and \emph{when} ---
the ordering rules, timing rules and formulas that govern their use are deliberately left out of
it, because those are the parts being examined. Misspelling a bit costs nothing; reaching for the
wrong register, or the right one at the wrong moment, costs everything.

\textbf{Code is marked on semantics, not on syntax.} A missing semicolon, a forgotten
\code{#include} or a brace in the wrong column costs nothing. A missing \code{volatile} costs
everything.

\textbf{Where a question says ``in the book's style''}, it means the conventions used in the
chapters:
\begin{itemize}
  \item \code{snake_case} for variables, functions and constants; \code{UPPER_CASE} for macros and
    enum constants.
  \item \code{static} on anything that is not part of a module's public API.
  \item A header guard named after the file in \code{UPPER_CASE}, \code{.} replaced by \code{_},
    trailing \code{_} (\code{driver/gpio.h} $\rightarrow$ \code{DRIVER_GPIO_H_}).
  \item A Doxygen comment in the imperative mood above every non-trivial function.
  \item \code{self} as the first parameter of any function operating on a driver struct, always a
    pointer, and \code{const} when the function only reads through it.
  \item Yoda notation on comparisons against a constant, and braces on every \code{if}, even
    one-liners.
  \item \code{const} on pass-by-value parameters in a function \emph{definition}, dropped from the
    prototype.
  \item Fixed-width types from \code{<stdint.h>}, and \code{U} suffixes on unsigned literals.
  \item \code{volatile} on anything shared between an ISR and the rest of the program.
\end{itemize}

\textbf{Where a question asks you to write code}, Doxygen comments are marked only where the
question says so. Everywhere else, write the code.

\textbf{Where a question says ``state the consequence''}, naming the defect earns half the marks and
saying what it does to the running program earns the other half. ``This is wrong'' scores nothing.

\subsection*{Marks}

\begin{narrowtable}{@{}lcccccccc@{}}
  \thead{Question} & 1 & 2 & 3 & 4 & 5 & 6 & 7 & 8 \\
  \midrule
  \thead{Marks} & 12 & 12 & 12 & 13 & 13 & 13 & 12 & 13 \\
\end{narrowtable}

% ------------------------------------------------------------------------------------------
\question{Saying what you mean}{12}
\label{exb:q:1}

\qpart{a} A blink module supports four patterns: off, slow, fast, and a two-pulse heartbeat.

Write, in the book's style, an enumeration \code{led_pattern_t} naming all four, and a function

\begin{ccode}
bool led_pattern_valid(led_pattern_t pattern);
\end{ccode}

\noindent that returns \code{true} only for a named member. State three concrete advantages an
\code{enum} has over a block of \code{#define}s for this, and then state the one property a
\code{#define} has that a \code{const uint8_t} does not, which is why this book still writes its pin
numbers as macros. \worth{4}

\qpart{b} Write a function-like macro \code{TICKS_NEEDED(timeout_ms, tick_ms)} that yields the
number of ticks in a timeout. State the two parenthesisation rules it must obey, showing for one of
them a call that goes wrong if the rule is broken. Then name the third macro hazard, which
parentheses cannot fix, and state whether your macro is exposed to it, with the reason. \worth{3}

\qpart{c} A host-side test harness grows a buffer as data arrives.

\begin{ccode}
typedef struct
{
    uint8_t* data;
    size_t size;
} buffer_t;
\end{ccode}

\noindent Write \code{buffer_resize(buffer_t* self, size_t size)} returning \code{0} on success and
\code{1} on failure. It must leave \code{self} completely untouched if the allocation fails, and
must handle a requested size of zero. State the specific mistake that a naive one-line
\code{self->data = realloc(self->data, size);} makes, and what it costs. \worth{3}

\qpart{d} A driver is split into \file{driver/serial.h} and \file{driver/serial.c}. Give the header
guard for that header, in the book's convention, and say what it is for. Then state the ordering
convention this book uses inside a \file{.c} file for its public functions and its \code{static}
helpers, and the reason. \worth{2}

% ------------------------------------------------------------------------------------------
\question{Straight to the registers}{12}
\label{exb:q:2}

\qpart{a} Three LEDs are wired to D8, D9 and D10, and a button to D12 through an external
pull-down, so the pin reads high while the button is held.

Write a complete single-file program, in the book's style, that lights exactly one LED at a time in
the order D8 $\rightarrow$ D9 $\rightarrow$ D10, advancing every 150\,ms, for as long as the button
is held, and turns all three off the moment it is released. Use macros for the pins and for the GPIO
operations, and a \code{delay_ms()} helper. \worth{5}

\qpart{b} Give the single C expression for each of the following on a register \code{reg} and a bit
position \code{n}: set the bit, clear the bit, test the bit, toggle it via \code{PORTx}, and toggle
it atomically. Then state the difference between \code{DDRB = (1U << 3U);} and
\code{DDRB |= (1U << 3U);}, and give one situation where each is the right choice. \worth{4}

\qpart{c} \code{PORTx} means one thing on an output pin and something else entirely on an input
pin. State both, and state which register must be read to get a pin's actual level.

The button circuit in \secref{c2:a:4} uses an external 10\,k$\Omega$ pull-down, yet the code still
writes \code{PORTB = (1U << BTN1);}. State what that write does, what the pin reads when idle, and
why the write is harmless in practice but not guaranteed to be. \worth{3}

% ------------------------------------------------------------------------------------------
\question{Reacting instead of polling}{12}
\label{exb:q:3}

\qpart{a} A button is wired to D3 through an external pull-down, so a press drives the pin high. An
LED is on D10. Every press-down must toggle the LED, and the toggle must be immune to another
interrupt changing a different bit of the same port at the same moment.

Write \code{setup()} and the interrupt handler, in the book's style, giving every register write and
the exact vector name. State which edge you selected and why. \worth{5}

\qpart{b} The button is moved from D3 to A2. State everything that changes: the enable bit, the
mask register, the bit written into it, and the vector. Then state \textbf{two} things the handler
must now do that the D3 version did not, and why each is forced on you rather than optional.
\worth{4}

\qpart{c} State the rule this book applies to decide what belongs inside an ISR and what belongs in
\code{main()}'s loop. Name two specific calls from the book's own drivers that must never appear in
an ISR, and give the concrete cost of each. Then write the three-line pattern an ISR and
\code{main()} use to communicate instead, and say what \code{volatile} is doing in it. \worth{3}

% ------------------------------------------------------------------------------------------
\question{Reading the world, and keeping time}{13}
\label{exb:q:4}

\qpart{a} Write \code{adc_read(uint8_t channel)} in full, in the book's style, returning a 10-bit
result. For each register you write, state in one sentence what that write is for. State what the
÷128 prescaler is chosen to achieve and what goes wrong without it. \worth{4}

\qpart{b} A 16\,MHz ATmega328P is configured with a ÷1024 prescaler, and a 250\,ms periodic event is
required.
\begin{enumerate}
  \item Give the tick period and the Timer~0 overflow period.
  \item Give the \code{OCR1A} value for a 250\,ms interrupt on Timer~1 in CTC mode, showing the
    arithmetic, and write the three register assignments plus the vector name.
  \item Give the number of Timer~0 overflows the same 250\,ms needs in Normal Mode, the actual
    period that produces, and the error as a percentage.
  \item State why Timer~0 cannot reach 250\,ms without counting overflows in software at any
    prescaler available to it.
\end{enumerate}
\noindent\worth{5}

\qpart{c} A potentiometer on A0 sets the brightness of an LED on D9 by software PWM over a 10\,ms
period.

Write \code{run_pwm_cycle()}, in the book's style, computing the on-time and off-time from a fresh
ADC reading. State why the book adds \code{+ 0.5F} before the cast, what this whole approach costs
the rest of the program, and what replaces it. \worth{4}

% ------------------------------------------------------------------------------------------
\question{A channel to the host}{13}
\label{exb:q:5}

\qpart{a} Write \file{driver/serial.h} in full, in the book's style: header guard, includes, and a
documented prototype for each of \code{serial_init()}, \code{serial_print()} and
\code{serial_write()} (the buffer-and-length form). Doxygen comments are marked here.

Then state which functions of the driver belong in the header and which must be \code{static} in
\file{serial.c}, with the reason, and state why a \code{const} on a by-value parameter is dropped
from the prototype but kept in the definition. \worth{5}

\qpart{b} Write \code{serial_write_byte()} and \code{serial_print()}. State what \code{UDRE0}
signals, and why the wait on it comes before the write to \code{UDR0} rather than after. Then state
what \code{sprintf()} is for in this context, and the one caveat \secref{c6:a:5} attaches to it on
AVR. \worth{4}

\qpart{c} A program must report an ADC reading over UART once a second while still responding to a
debounced button on a pin-change interrupt.

Describe where the \code{sprintf()} and the \code{serial_print()} calls go, and why they cannot go
where the event is detected. Write the flag pattern that connects the two, and state exactly what
\code{volatile} is doing. Then state why the \code{read()} of \secref{c6:a:7} takes a
\code{timeout_ms} parameter, and what each of its three possible return values means. \worth{4}

% ------------------------------------------------------------------------------------------
\question{Surviving a power cut, and a hang}{13}
\label{exb:q:6}

\qpart{a} Write \code{eeprom_write_byte(uint16_t addr, uint8_t byte)}\newline and
\code{eeprom_read_byte(uint16_t addr)} in full, in the book's style. For the write, explain each of
the following in one sentence: the loop on \code{EEPE} at the top, the order of \code{EEMPE} and
\code{EEPE}, and the \code{cli()}/\code{sei()} pair. State why the read needs no equivalent of the
last of those. \worth{5}

\qpart{b} Write the enumeration \code{watchdog_timeout_t} covering all ten timeout settings, and
\code{watchdog_init(watchdog_timeout_t timeout)}.

Explain why eight of the ten members can take their default values while the last two must be given
\code{32} and \code{33} explicitly. Then explain what the \code{WDCE} write immediately before the
real one is for, and what happens if the second write comes too late. \worth{4}

\qpart{c} Four values are candidates for EEPROM storage:
\begin{enumerate}
  \item a calibration offset written once during manufacture;
  \item a live ADC reading sampled at 100\,Hz;
  \item an on/off state a user toggles with a button;
  \item a boot counter incremented on every reset.
\end{enumerate}
For each, state whether EEPROM is appropriate and why, quantifying your answer for the one that is
clearly wrong. Then state where in \code{main()} the watchdog reset belongs, and how the timeout
should be chosen relative to the loop it is protecting. \worth{4}

% ------------------------------------------------------------------------------------------
\question{A driver that scales}{12}
\label{exb:q:7}

\qpart{a} Write \file{driver/gpio.h}, and the body of \code{gpio_init()}, for a GPIO driver in which
the caller passes a single Arduino pin number rather than register pointers:

\begin{ccode}
bool gpio_init(gpio_t* self, uint8_t pin, gpio_mode_t mode);
\end{ccode}

\begin{itemize}
  \item Pins 0--7 map to \code{PORTD0}--\code{PORTD7}; pins 8--13 map to
    \code{PORTB0}--\code{PORTB5}.
  \item \code{gpio_init()} derives and stores the register pointers and the bit position itself.
  \item Three modes are supported: input, input with pull-up, and output.
  \item Any pin outside 0--13 returns \code{false} without touching a single register.
\end{itemize}
Doxygen comments are marked on the header. \worth{6}

\qpart{b} In the original driver, \code{gpio_t}'s \code{pin} field was declared
\code{const uint8_t pin;}. State why it can no longer be \code{const} in your version, and what
specifically a \code{const} struct member permits and forbids. Then state why \code{gpio_read()}
takes \code{const gpio_t* self} while \code{gpio_write()} takes \code{gpio_t* self}, and what would
change if either took \code{gpio_t self} by value instead. \worth{3}

\qpart{c} Two calls to \code{gpio_init()} with the same pin number would silently hand out two
drivers for one piece of hardware.

Write the pin-reservation table that prevents it: the declaration in \file{gpio.c}, the check and
the claim inside \code{gpio_init()}, and the release inside \code{gpio_free()}. State why the table
must be \code{static}, and what breaks if it is not. \worth{3}

% ------------------------------------------------------------------------------------------
\question{Hiding the implementation}{13}
\label{exb:q:8}

\qpart{a} The GPIO driver's struct currently sits, fields and all, in \file{gpio.h}.

Write the opaque version of that header: what replaces the full struct definition, and what has to
happen to the struct itself. Then state the three things this takes away from calling code, and for
each, what the caller must do instead. \worth{4}

\qpart{b} Write a \code{button_t} struct holding a GPIO and a press callback, together with

\begin{ccode}
static void button_check(const button_t* self, bool edge);
\end{ccode}

\noindent which invokes the callback when a press has been detected. State in words what
\code{void (*on_press)(void);} declares, write the line that assigns a function to it, and write the
line that calls through it. Then state why \code{button_check()} is better for never mentioning the
LED it ultimately toggles. \worth{4}

\qpart{c} Given the interface

\begin{ccode}
typedef struct gpio_interface gpio_interface_t;

typedef struct gpio_vtable
{
    void (*del)(gpio_interface_t** self);
    bool (*read)(const gpio_interface_t* self);
    void (*write)(gpio_interface_t* self, bool value);
    void (*toggle)(gpio_interface_t* self);
} gpio_vtable_t;

typedef struct gpio_interface
{
    const gpio_vtable_t* vptr;
} gpio_interface_t;
\end{ccode}

\noindent write a second implementation, \code{gpio_stub_t}, that touches no hardware at all: the
struct, its four vtable functions, the vtable itself, and \code{gpio_stub_new()}.

State the one rule your struct's field order must obey and the language guarantee that makes the
whole scheme work. State why a single shared vtable is enough for every instance ever created. Then
state what having a stub makes possible that a hardware-only driver does not. \worth{5}

% ------------------------------------------------------------------------------------------
\section*{Register Reference}
\phantomsection\label{exb:sec:registers}
\addcontentsline{toc}{section}{Register Reference}
\markright{Register Reference}

Every ATmega328P register and bit this paper refers to is listed here, with what it is. Nothing
below is worth marks by itself: the paper tests which register to reach for and when, not whether
you can recall its name. Where a bit's \textbf{position} matters it is given; where a timing rule,
an ordering rule or a formula matters it is not, because that is the part being examined.

The target is an Arduino Uno/Nano at 16\,MHz.

\subsection*{Board pins and ports}

\begin{narrowtable}{@{}ll@{}}
  \thead{Board pin} & \thead{Port bits} \\
  \midrule
  D0--D7 & \code{PORTD0}--\code{PORTD7} \\
  D8--D13 & \code{PORTB0}--\code{PORTB5} \\
  A0--A5 & \code{PORTC0}--\code{PORTC5} \\
\end{narrowtable}

\subsection*{Digital I/O --- one set per port, \code{x} being \code{B}, \code{C} or \code{D}}

\begin{booktable}{@{}lL@{}}
  \thead{Register} & \thead{What it is} \\
  \midrule
  \code{DDRx} & Data direction, one bit per pin: \code{1} = output, \code{0} = input. \\
  \code{PORTx} & On an output pin, the level driven onto it. On an input pin, the internal pull-up
    enable. \\
  \code{PINx} & Reads the actual logic level on each pin. Writing a \code{1} to a bit toggles the
    matching \code{PORTx} bit, whatever \code{DDRx} says; bits written \code{0} are unaffected. \\
\end{booktable}

\subsection*{Interrupts}

\begin{booktable}{@{}>{\raggedright\arraybackslash}p{8.6em}L@{}}
  \thead{Register / name} & \thead{What it is} \\
  \midrule
  \code{SREG} & Status register. Bit \code{I} is the global interrupt enable; \code{sei()} sets it,
    \code{cli()} clears it. \\
  \code{EICRA} & External interrupt edge select: \code{ISC01:ISC00} for INT0,
    \code{ISC11:ISC10} for INT1. \\
  \code{EIMSK} & External interrupt enable: bits \code{INT0}, \code{INT1}. \\
  \code{INT0_vect} /\newline \code{INT1_vect} & External interrupt vectors, for D2 and D3
    respectively. \\
  \code{PCICR} & Pin change interrupt enable per port: \code{PCIE0} (B), \code{PCIE1} (C),
    \code{PCIE2} (D). \\
  \code{PCMSK0} / \code{PCMSK1} /\newline \code{PCMSK2} & Pin change mask for ports B, C and D.
    The bit position matches the pin's own bit position within its port. \\
  \code{PCINT0_vect} /\newline \code{PCINT1_vect} /\newline \code{PCINT2_vect} & Pin change
    interrupt vectors, for ports B, C and D. \\
\end{booktable}

\noindent\code{ISCn1:ISCn0} encoding: \code{00} low level, \code{01} any logical change, \code{10}
falling edge, \code{11} rising edge.

\subsection*{ADC}

\begin{booktable}{@{}>{\raggedright\arraybackslash}p{8.6em}L@{}}
  \thead{Register / name} & \thead{What it is} \\
  \midrule
  \code{ADMUX} & Bits 3:0 select the input channel. \code{REFS1:REFS0} (bits 7:6) select the
    reference: \code{00} = the external \code{AREF} pin, \code{01} = \code{AVCC}. \\
  \code{ADCSRA} & \code{ADEN} enables the ADC, \code{ADSC} starts a conversion, \code{ADIF} is the
    conversion-complete flag (cleared by writing a \code{1} to it), \code{ADPS2:ADPS0} set the
    prescaler. \\
  \code{ADCL} / \code{ADCH} & Low and high halves of the 10-bit result. \\
  \code{ADC} & Macro combining \code{ADCL} and \code{ADCH} into the full 10-bit value. \\
\end{booktable}

\noindent\code{ADPS2:ADPS0} encoding: \code{000} and \code{001} both divide by 2, and each value
above that doubles the divisor, up to \code{111} = ÷128.

\subsection*{Timers --- \code{n} is \code{0}, \code{1} or \code{2}}

\begin{booktable}{@{}>{\raggedright\arraybackslash}p{8.6em}L@{}}
  \thead{Register / name} & \thead{What it is} \\
  \midrule
  \code{TCCRnB} & Clock select \code{CSn2:CSn0} (the prescaler). \code{WGM12}, in \code{TCCR1B},
    selects CTC mode on Timer~1. \\
  \code{TCNTn} & The counter itself. \\
  \code{OCRnA} & Compare value. In CTC mode the timer counts from 0 up to it, then clears. \\
  \code{TIMSKn} & \code{TOIEn} enables the overflow interrupt; \code{OCIEnA} enables the
    compare-match interrupt. \\
  \code{TIMERn_OVF_vect} & Overflow vector. \\
  \code{TIMERn_COMPA_vect} & Compare-match vector. \\
\end{booktable}

\noindent Timer~0 and Timer~2 are 8-bit; Timer~1 is 16-bit.

\code{CSn2:CSn0} encoding differs between timers:

\begin{narrowtable}{@{}lll@{}}
  \thead{Bits} & \thead{Timer 0 / Timer 1} & \thead{Timer 2} \\
  \midrule
  \code{001} & ÷1 & ÷1 \\
  \code{010} & ÷8 & ÷8 \\
  \code{011} & ÷64 & ÷32 \\
  \code{100} & ÷256 & ÷64 \\
  \code{101} & ÷1024 & ÷128 \\
  \code{110} & -- & ÷256 \\
  \code{111} & -- & ÷1024 \\
\end{narrowtable}

\subsection*{UART}

\begin{booktable}{@{}>{\raggedright\arraybackslash}p{8.6em}L@{}}
  \thead{Register / name} & \thead{What it is} \\
  \midrule
  \code{UCSR0A} & \code{UDRE0} is set when \code{UDR0} is free to accept a new byte. \code{RXC0}
    is set when a fully received byte is waiting in \code{UDR0}. \\
  \code{UCSR0B} & \code{TXEN0} enables the transmitter, \code{RXEN0} enables the receiver. \\
  \code{UCSR0C} & Frame format. \code{UCSZ01} and \code{UCSZ00} set together give 8 data bits. \\
  \code{UBRR0} & Baud rate divisor. Only the low 12 bits are used. \\
  \code{UDR0} & Data register: written to transmit, read to receive. Reading it clears
    \code{RXC0}. \\
\end{booktable}

\subsection*{EEPROM}

\begin{booktable}{@{}>{\raggedright\arraybackslash}p{8.6em}L@{}}
  \thead{Register / name} & \thead{What it is} \\
  \midrule
  \code{EEAR} & Address to read from or write to, 0--1023. \\
  \code{EEDR} & Data register: holds the byte to be written, or receives the byte just read. \\
  \code{EECR} & \code{EERE} starts a read. \code{EEMPE} is the master write enable, a safety
    interlock rather than a data bit. \code{EEPE} starts a write and stays set until that write
    completes. \\
\end{booktable}

\noindent The ATmega328P has 1024 bytes of EEPROM. A byte write takes a few milliseconds, and a
cell tolerates roughly 100,000 write cycles.

\subsection*{Watchdog}

\begin{booktable}{@{}>{\raggedright\arraybackslash}p{8.6em}L@{}}
  \thead{Register / name} & \thead{What it is} \\
  \midrule
  \code{MCUSR} & MCU status register. \code{WDRF} is set by a watchdog-caused reset. \\
  \code{WDTCSR} & Watchdog control and status. Bit 0 \code{WDP0}, bit 1 \code{WDP1}, bit 2
    \code{WDP2}, bit 3 \code{WDE}, bit 4 \code{WDCE}, bit 5 \code{WDP3}, bit 6 \code{WDIE}, bit 7
    \code{WDIF}. \\
  \code{WDT_vect} & Watchdog interrupt vector. \\
\end{booktable}

\noindent\code{WDE} is the watchdog enable, \code{WDIE} the watchdog interrupt enable, and
\code{WDCE} the change enable used when reconfiguring them. \code{WDP3:WDP0} select the timeout:
\code{0000} is 16\,ms, and each step doubles it, up to \code{1001} = 8192\,ms.
